\section{Practical Implications and Case Study}
\label{sec:case}

\subsection{Practical implications}
Three implications follow for anyone who builds or buys an Arabic poetry generator. First, \emph{a classifier score is
not a quality gate}: a system that reports 95\% meter accuracy may deliver fewer than 10\% fully correct poems. Second,
\emph{validity must be measured per poem and jointly}: meter, rhyme and length fail together and compound with length,
so per-metric averages hide most failures. Third, \emph{a rule-based verifier is cheap to run and explains its
verdicts}: it needs no GPU, it names the broken foot, and it can be placed in front of human editors so that their time
goes to meaning rather than to scansion.

\subsection{Case study: a textbook publisher}
To make these points concrete we walk through an \emph{illustrative} scenario. The organisation is a composite, not a
real company, and the numbers are taken from our measurements.

\textbf{Setting.} An educational publisher prepares a secondary-school Arabic literature textbook. The prosody chapter
needs 200 short original poems (two or three bayts each) in six common meters, on topics from the curriculum, for
students to scan as exercises. A poem with a broken hemistich is worse than no poem: students would learn a wrong
pattern, and teachers would lose trust in the book. The publisher has two editors who know prosody and a small budget
of cloud GPU time.

\textbf{Option A: the generator and its published metric.} The team generates one poem per request and keeps the poems
that pass the base criterion. By our measurements, \val{vEzeroBaseBilstm}\% of the released poems pass it, so
the team quickly has 200 ``validated'' poems. But only \val{gEzeroStrictGivenBase}\% of poems that pass this criterion are
strictly valid. The editors must therefore scan every poem by hand and will reject most of them, often for a rhyme
fault that no metric flagged. The saving from automation disappears.

\textbf{Option B: a verifier gate with resampling.} The team puts the verifier between the generator and the editors
(Fig.~\ref{fig:case}). A poem reaches an editor only if it is strictly valid; otherwise the request is sampled again,
up to a budget. For 1--3-bayt poems the released model is strictly valid \val{mEzeroLenA}\% of the time, so a valid poem
needs about six samples on average if samples are independent (an optimistic assumption, since some prompts are harder
than others). At the 2.4\,s per poem we measured for generation on an L4 GPU, 200 valid poems cost roughly
$200 \times 6 \times 2.4\,\mathrm{s} \approx 48$ minutes of GPU time. The editors now review only meaning, tone and
suitability for students. When a request exhausts its budget, the verifier's report (the hemistich, the nearest valid
pattern and the broken foot) lets an editor repair the line by hand in minutes.

\textbf{Option C: hemistich search (pending).} With the M4 search, the generator repairs each hemistich as it writes,
so the predicted cost per valid poem is lower than in Option B, and poems of four or more bayts, for which Option B
would need about 45 samples each, become feasible. The measured cost is \pending{pending M4}.

\begin{figure}[t]
\centering
\begin{tikzpicture}[
  font=\scriptsize,
  st/.style={draw=black!70, rounded corners=2pt, fill=#1, align=center, text width=1.75cm, minimum height=0.95cm, inner sep=2pt},
  st/.default={blue!6},
  dec/.style={draw=black!70, diamond, aspect=1.6, fill=yellow!12, align=center, inner sep=0.5pt, font=\scriptsize},
  arr/.style={-{Stealth[length=2mm]}, thick, draw=black!70}]
\node[st={gray!10}] (req) {Request: meter, length, topic, raw\={\i}};
\node[st, right=0.35cm of req] (gen) {Generator (Shaer adapter)};
\node[dec, right=0.35cm of gen] (ver) {Verifier\\ valid?};
\node[st={green!8}, right=0.35cm of ver] (ed) {Editor: meaning, tone};
\node[st={green!8}, below=0.6cm of ed] (book) {Textbook};
\node[dec, below=0.7cm of gen] (bud) {Budget\\ left?};
\node[st={orange!10}, left=0.35cm of bud] (fix) {Editor repairs from report};
\draw[arr] (req) -- (gen);
\draw[arr] (gen) -- (ver);
\draw[arr] (ver) -- node[above, font=\scriptsize] {yes} (ed);
\draw[arr] (ed) -- (book);
\draw[arr] (ver) |- node[pos=0.25, right, font=\scriptsize] {no} (bud);
\draw[arr] (bud) -- node[right, font=\scriptsize] {yes} (gen);
\draw[arr] (bud) -- node[above, font=\scriptsize] {no} (fix);
\draw[arr] (fix.south) |- ($(book.south west)+(0,0.2)$);
\end{tikzpicture}
\caption{Case-study scenario: the verifier as a quality gate between the generator and the human editors.}
\label{fig:case}
\end{figure}

\textbf{Lessons.} The publisher's decision depends on the metric it trusts. Under the published metric the generator
looks ready for use; under the strict criterion it is a useful drafting tool that needs a verifier in the loop. The
gate also changes the editors' work from scanning lines to judging content, which is where human judgement is most
valuable.
